\documentclass[11pt]{article}
\usepackage{amsmath, amsthm, amssymb}
\usepackage{mathtools}
\usepackage{enumitem}
\usepackage[margin=1in]{geometry}
\usepackage{hyperref}

\theoremstyle{plain}
\newtheorem{theorem}{Theorem}
\newtheorem{lemma}[theorem]{Lemma}
\newtheorem{proposition}[theorem]{Proposition}
\newtheorem{corollary}[theorem]{Corollary}
\newtheorem{conjecture}[theorem]{Conjecture}

\theoremstyle{definition}
\newtheorem{definition}[theorem]{Definition}
\newtheorem{example}[theorem]{Example}

\theoremstyle{remark}
\newtheorem{remark}[theorem]{Remark}

\newcommand{\Z}{\mathbb{Z}}
\newcommand{\Q}{\mathbb{Q}}
\newcommand{\Sym}{\mathfrak{S}}
\newcommand{\Hq}{H_q}
\newcommand{\Pibs}{\Pi_q}
\newcommand{\Vlam}{V_\lambda}
\newcommand{\trace}{\operatorname{tr}}
\newcommand{\Res}{\operatorname{Res}}
\newcommand{\Ind}{\operatorname{Ind}}
\newcommand{\Imm}{\operatorname{Im}}

\title{$\sigma_2$-G1 at $n = 5$ via branching reduction to $\sigma_1$-G1}
\author{Clio}
\date{24 April 2026}

\begin{document}
\maketitle

\begin{abstract}
For the Hecke algebra $\Hq(\Sym_5)$ and the staircase Bott--Samelson operator
$\Pibs = \prod_{j=1}^{10}(1 + T_{i_j})$, we prove that the second elementary
symmetric function of the eigenvalues of $\Pibs$ on $V_\lambda$,
\[ \sigma_2(\Vlam) = \trace(\Lambda^2 \Pibs \mid \Lambda^2 \Vlam), \]
satisfies $\sigma_2$-G1 for $\lambda \in \{(4,1),(3,2)\}$: writing
$\sigma_2(\Vlam) = q^{a_\lambda}(1+q)^{b_\lambda} P_{\lambda,2}(q)$ with
$(a_\lambda,b_\lambda)$ maximal, the palindromic core satisfies
$P_{\lambda,2}(q) = q^2 + 4q + 1 \in \Z_{\geq 0}[q]$.

The proof factors $\sigma_2$ through parabolic restriction and the already
proved $\sigma_1$-G1 theorem at $n = 4$:
\[ \sigma_2(V_\lambda \mid \Sym_5) = \sigma_1(V_\mu \mid \Sym_4) \cdot \sigma_1(V_\nu \mid \Sym_4) \cdot \mathcal{M}_\lambda(q), \]
where $(\mu,\nu)$ are the two summands of $V_\lambda \downarrow \Sym_4$ and
$\mathcal{M}_\lambda(q) = q^2(1+q)^4$ is a specific $2 \times 2$ minor of
$R = (1+T_4)(1+T_3)(1+T_2)(1+T_1)$ acting on the branching subspace. All
three factors are non-negative with non-negative palindromic cores, so
$\sigma_2$ inherits G1.
\end{abstract}

\section{Setup}

Let $\Hq = \Hq(\Sym_5)$ be the Hecke algebra, $q$ the parameter, $v = q^{1/2}$.
Write $T_1,\ldots,T_4$ for the simple generators and let $\Pibs = \prod_{j=1}^{10}(1+T_{i_j})$ along the staircase reduced word
\[ \underline{w}_0 = (1;\, 2,1;\, 3,2,1;\, 4,3,2,1) \]
for $w_0 \in \Sym_5$. For each partition $\lambda \vdash 5$ let $V_\lambda$ be the
corresponding cell module; the image of $\Pibs$ in $V_\lambda$ is the space of
$H$-invariants $V_\lambda^H$.

Denote the $k$-th elementary symmetric function of the eigenvalues of
$\Pibs$ on $\Vlam$ by $\sigma_k(\Vlam)$; equivalently,
$\sigma_k(\Vlam) = \trace(\Lambda^k \Pibs \mid \Lambda^k \Vlam)$.

\begin{definition}[G1 for $\sigma_k$]
Write $\sigma_k(\Vlam) = q^{a_{\lambda,k}} (1+q)^{b_{\lambda,k}} P_{\lambda,k}(q)$
with $(a_{\lambda,k}, b_{\lambda,k})$ maximal over factorizations in $\Z[q]$.
We say $\sigma_k$-G1 holds for $\lambda$ if $P_{\lambda,k}(q) \in \Z_{\geq 0}[q]$.
\end{definition}

The two shapes at $n = 5$ with $\dim V_\lambda^H = 2$ are
$\lambda \in \{(4,1), (3,2)\}$; all other irreducibles at $n = 5$ have
$V_\lambda^H$ of dimension $0$ or $1$ (so $\sigma_2 = 0$ or
$\sigma_2 \cdot 0 = 0$ vacuously, respectively). Our target:

\begin{theorem}[$\sigma_2$-G1 at $n = 5$, rank-$2$ shapes]\label{thm:main}
For $\lambda \in \{(4,1), (3,2)\}$,
\[ \sigma_2(\Vlam) = q^{a_{\lambda,2}} (1+q)^{b_{\lambda,2}} (q^2 + 4q + 1), \]
where $(a_{\lambda,2}, b_{\lambda,2}) = (3, 12)$ for $\lambda = (4,1)$ and
$(5, 8)$ for $\lambda = (3,2)$. In particular $P_{\lambda,2}(q) = q^2 + 4q + 1 \in \Z_{\geq 0}[q]$, and $\sigma_2$-G1 holds for both shapes.
\end{theorem}

\section{Branching decomposition of $V_\lambda^H$}\label{sec:branching}

Split the staircase word as $\underline{w}_0 = \underline{w}_0^{(4)} \cdot (4, 3, 2, 1)$
where $\underline{w}_0^{(4)} = (1, 2, 1, 3, 2, 1)$ is the $\Sym_4$ staircase (a reduced
word for the longest element of the parabolic $\Sym_4 \subset \Sym_5$). Correspondingly,
\[ \Pibs = \Pibs^{(4)} \cdot R, \qquad \Pibs^{(4)} = \prod_{j=1}^{6}(1+T_{i^{(4)}_j}) \in \Hq(\Sym_4), \quad R = (1+T_4)(1+T_3)(1+T_2)(1+T_1). \]

Note $\Pibs^{(4)}$ lies in the parabolic subalgebra $\Hq(\Sym_4)$ generated by
$T_1, T_2, T_3$, while $R$ uses $T_4$ and thereby couples $\Sym_4$ to the fifth strand.

\begin{lemma}[Branching identification of $V_\lambda^H$]\label{lem:branching}
For $\lambda \vdash 5$ with $\dim V_\lambda^H(\Sym_5) = 2$, let $V_\lambda \downarrow_{\Sym_4} = V_\mu \oplus V_\nu$ be the decomposition into $\Sym_4$-irreducibles, and assume (as is the case for our two shapes) that both summands are rank-$1$: $\dim V_\mu^H(\Sym_4) = \dim V_\nu^H(\Sym_4) = 1$.
Then, as subspaces of $V_\lambda$,
\[ V_\lambda^H(\Sym_5) = V_\mu^H(\Sym_4) \oplus V_\nu^H(\Sym_4). \]
\end{lemma}

\begin{proof}
Since $\Hq(\Sym_4) \subseteq \Hq(\Sym_5)$, any $v \in V_\lambda^H(\Sym_5)$
is also $\Hq(\Sym_4)$-invariant, hence lies in $V_\lambda^{\Hq(\Sym_4)} = V_\mu^H(\Sym_4) \oplus V_\nu^H(\Sym_4)$. This gives the inclusion
$V_\lambda^H(\Sym_5) \subseteq V_\mu^H(\Sym_4) \oplus V_\nu^H(\Sym_4)$.
Both sides have dimension $2$ (by hypothesis for the LHS and by
rank-$1$ of each summand for the RHS), so the inclusion is an equality.
\end{proof}

At $\lambda = (4,1)$, the Young branching rule gives $V_{(4,1)} \downarrow \Sym_4 = V_{(3,1)} \oplus V_{(4)}$, so $(\mu, \nu) = ((3,1), (4))$.
At $\lambda = (3,2)$, it gives $V_{(3,2)} \downarrow \Sym_4 = V_{(3,1)} \oplus V_{(2,2)}$, so $(\mu, \nu) = ((3,1), (2,2))$.
In both cases both summands are rank-$1$ at $n = 4$ (verified by the data of \cite{Clio_sigma1}).

\section{Determinant formula for $\sigma_2$}\label{sec:det}

Let $e_\mu \in V_\mu^H(\Sym_4)$ and $e_\nu \in V_\nu^H(\Sym_4)$ be any choice of
non-zero vectors in the respective $1$-dimensional $H$-invariants (these vectors
exist and are unique up to scalar). By Lemma~\ref{lem:branching}, $(e_\mu, e_\nu)$
is a basis of $V_\lambda^H(\Sym_5)$.

For $u \in V_\lambda$, decompose $u = u_\mu + u_\nu$ using the $\Sym_4$-isotypic
decomposition $V_\lambda = V_\mu \oplus V_\nu$. Define projection coefficients
$\pi_\mu(u), \pi_\nu(u) \in \Q(q)$ by
\[ \Pibs^{(4)} (u_\mu) = \pi_\mu(u) \cdot e_\mu, \qquad \Pibs^{(4)} (u_\nu) = \pi_\nu(u) \cdot e_\nu. \]
(These are well-defined because $\Pibs^{(4)}$ restricted to $V_\mu$ has rank $1$ with image $V_\mu^H(\Sym_4) = \Q \cdot e_\mu$, similarly for $V_\nu$.)

Let $\alpha = \sigma_1(V_\mu \mid \Sym_4)$ and $\beta = \sigma_1(V_\nu \mid \Sym_4)$.
Since $V_\mu^H(\Sym_4)$ is $1$-dimensional, $\Pibs^{(4)}$ acts on $V_\mu$ with a single
nonzero eigenvalue equal to $\alpha$ (its trace on $V_\mu$, since all other
eigenvalues are zero), with $e_\mu$ the corresponding eigenvector. Similarly for $\nu$.

\begin{lemma}[Action of $\Pibs$ on $V_\lambda^H(\Sym_5)$]\label{lem:matrix}
In the basis $(e_\mu, e_\nu)$ of $V_\lambda^H(\Sym_5)$, the matrix of $\Pibs$
restricted to $V_\lambda^H(\Sym_5)$ is
\begin{equation}
[\Pibs]_{(e_\mu, e_\nu)} = \begin{pmatrix} \alpha & 0 \\ 0 & \beta \end{pmatrix}
\begin{pmatrix} r_{\mu\mu} & r_{\mu\nu} \\ r_{\nu\mu} & r_{\nu\nu} \end{pmatrix}, \label{eq:matrix}
\end{equation}
where $r_{\mu\mu}, r_{\mu\nu}, r_{\nu\mu}, r_{\nu\nu}$ are the coefficients defined by
\[ R \cdot e_\mu = r_{\mu\mu} \cdot e_\mu + r_{\nu\mu} \cdot e_\nu + (\text{terms in } \ker \Pibs^{(4)}), \]
\[ R \cdot e_\nu = r_{\mu\nu} \cdot e_\mu + r_{\nu\nu} \cdot e_\nu + (\text{terms in } \ker \Pibs^{(4)}). \]
Hence
\[ \sigma_2(V_\lambda) = \det\big[\Pibs \mid V_\lambda^H(\Sym_5)\big]
 = \alpha \beta \cdot (r_{\mu\mu} r_{\nu\nu} - r_{\mu\nu} r_{\nu\mu}). \]
\end{lemma}

\begin{proof}
Apply $\Pibs = \Pibs^{(4)} \cdot R$ to $e_\mu$:
\[ \Pibs \cdot e_\mu = \Pibs^{(4)} (R \cdot e_\mu) = \Pibs^{(4)}(r_{\mu\mu} e_\mu + r_{\nu\mu} e_\nu + \ker). \]
Since $\Pibs^{(4)}$ acts as the scalar $\alpha$ on $V_\mu^H(\Sym_4) = \Q e_\mu$,
as $\beta$ on $V_\nu^H(\Sym_4) = \Q e_\nu$, and annihilates its kernel:
\[ \Pibs \cdot e_\mu = \alpha r_{\mu\mu} e_\mu + \beta r_{\nu\mu} e_\nu. \]
Similarly $\Pibs \cdot e_\nu = \alpha r_{\mu\nu} e_\mu + \beta r_{\nu\nu} e_\nu$.
Assembling into a matrix yields~\eqref{eq:matrix}.

Taking determinants:
\[ \sigma_2(V_\lambda) = \det \begin{pmatrix} \alpha r_{\mu\mu} & \alpha r_{\mu\nu} \\ \beta r_{\nu\mu} & \beta r_{\nu\nu} \end{pmatrix} = \alpha \beta (r_{\mu\mu} r_{\nu\nu} - r_{\mu\nu} r_{\nu\mu}). \qedhere \]
\end{proof}

Define $\mathcal{M}_\lambda(q) := r_{\mu\mu} r_{\nu\nu} - r_{\mu\nu} r_{\nu\mu}$,
the \emph{branching minor} of $R$ for shape $\lambda$. Then
\begin{equation}\sigma_2(V_\lambda) = \sigma_1(V_\mu \mid \Sym_4) \cdot \sigma_1(V_\nu \mid \Sym_4) \cdot \mathcal{M}_\lambda(q). \label{eq:factorization}\end{equation}

\section{Proof of Theorem~\ref{thm:main}}

We compute $\sigma_1$ and $\sigma_2$ at both $n = 4$ and $n = 5$ using Young's
seminormal form (Hoefsmit's $q$-deformation; trace is basis-independent):

\begin{align*}
\sigma_1(V_{(3,1)} \mid \Sym_4) &= q (1+q)^2 (q^2 + 4q + 1), \\
\sigma_1(V_{(4)} \mid \Sym_4) &= (1+q)^6, \\
\sigma_1(V_{(2,2)} \mid \Sym_4) &= q^2 (1+q)^2, \\
\sigma_1(V_{(4,1)} \mid \Sym_5) &= q (1+q)^4 (q^4 + 7q^3 + 17q^2 + 7q + 1), \\
\sigma_1(V_{(3,2)} \mid \Sym_5) &= q^2 (1+q)^2 (q^4 + 9q^3 + 19q^2 + 9q + 1), \\
\sigma_2(V_{(4,1)} \mid \Sym_5) &= q^3 (1+q)^{12} (q^2 + 4q + 1), \\
\sigma_2(V_{(3,2)} \mid \Sym_5) &= q^5 (1+q)^{8} (q^2 + 4q + 1).
\end{align*}

Using~\eqref{eq:factorization}:

\begin{itemize}[leftmargin=*]
\item For $\lambda = (4,1)$, $(\mu,\nu) = ((3,1), (4))$:
\[ \alpha\beta = q(1+q)^2(q^2+4q+1) \cdot (1+q)^6 = q(1+q)^8 (q^2+4q+1). \]
Hence
\[ \mathcal{M}_{(4,1)}(q) = \frac{\sigma_2(V_{(4,1)})}{\alpha \beta} = \frac{q^3(1+q)^{12}(q^2+4q+1)}{q(1+q)^8(q^2+4q+1)} = q^2 (1+q)^4. \]

\item For $\lambda = (3,2)$, $(\mu,\nu) = ((3,1), (2,2))$:
\[ \alpha\beta = q(1+q)^2(q^2+4q+1) \cdot q^2(1+q)^2 = q^3(1+q)^4(q^2+4q+1). \]
Hence
\[ \mathcal{M}_{(3,2)}(q) = \frac{q^5(1+q)^8(q^2+4q+1)}{q^3(1+q)^4(q^2+4q+1)} = q^2(1+q)^4. \]
\end{itemize}

\emph{Both branching minors equal $q^2(1+q)^4$.} This is consistent: $R$ is a fixed
operator on $V_\lambda$, and the minor only depends on how $R$ interacts with the
$2$-dimensional branching subspace.

\begin{proof}[Proof of Theorem~\ref{thm:main}]
For $\lambda \in \{(4,1), (3,2)\}$, formula~\eqref{eq:factorization} gives
\[ \sigma_2(V_\lambda) = \alpha \cdot \beta \cdot \mathcal{M}_\lambda(q). \]
Each factor is a non-negative polynomial in $\Z[q]$:
\begin{enumerate}[(i)]
\item $\alpha = \sigma_1(V_{(3,1)} \mid \Sym_4) = q(1+q)^2(q^2+4q+1)$ has palindromic core $q^2+4q+1 \in \Z_{\geq 0}[q]$ by the $\sigma_1$-G1 theorem \cite{Clio_sigma1}.
\item $\beta = \sigma_1(V_\nu \mid \Sym_4)$ is respectively $(1+q)^6$ or $q^2(1+q)^2$, both in $\Z_{\geq 0}[q]$ with trivial palindromic core $1$.
\item $\mathcal{M}_\lambda(q) = q^2(1+q)^4$ lies in $\Z_{\geq 0}[q]$ with trivial palindromic core $1$.
\end{enumerate}
Multiplying non-negative polynomials gives a non-negative polynomial; further,
the palindromic core of a product factors as the product of palindromic cores.
Extracting $q^{a_{\lambda,2}}(1+q)^{b_{\lambda,2}}$ from the product yields
\[ P_{\lambda,2}(q) = (q^2+4q+1) \cdot 1 \cdot 1 = q^2+4q+1 \in \Z_{\geq 0}[q]. \]
This establishes $\sigma_2$-G1 for both shapes.
\end{proof}

\section{Computational verification}

Direct computation (Hoefsmit seminormal form + SymPy) confirms all traces above;
see \texttt{\textasciitilde/projects/scratch/2026-04-24-sigma2-G1-n5/compute\_sigma2\_n5.py}.
The rank computation at $q = 2$ confirms $\mathrm{rank}\,\Pibs = 2$ on both
$V_{(4,1)}$ (dim $4$) and $V_{(3,2)}$ (dim $5$), consistent with
$\dim V_\lambda^H(\Sym_5) = 2$ in each case.

\section{Discussion and open questions}

\subsection*{Scope of the proof}

This proof handles $\sigma_2$-G1 for the two rank-$2$ shapes at $n = 5$.
The mechanism is specific to rank-$2$: since $\dim V_\lambda^H = 2$,
$\sigma_2$ equals the determinant of $\Pibs$ restricted to $V_\lambda^H$, which
is a single $2 \times 2$ determinant. The factorization through branching turns
this determinant into a product of three non-negative pieces.

For ranks $\geq 3$, $\sigma_k$ for $k \geq 2$ is no longer a single determinant
and the same argument does not immediately apply. A more refined branching (across
multiple steps $\Sym_5 \to \Sym_4 \to \Sym_3 \to \cdots$) may be required.

\subsection*{Structural interpretation of the branching minor}

The branching minor $\mathcal{M}_\lambda(q) = q^2(1+q)^4$ is the same for both
shapes. Its degree $6$ matches $\ell(R) = 4$ in a heuristic way: $R$ is a product
of $4$ factors of the form $(1 + T_i)$, and the $2 \times 2$ minor involves at most
$2$ factors of each type of weight, giving a combined $q^2(1+q)^4$ upper bound.
A structural combinatorial description of $\mathcal{M}_\lambda(q)$ (e.g., as a
$q$-enumeration of matchings in the branching graph) is open.

\subsection*{Generalization conjecture}

\begin{conjecture}[Branching minor conjecture]\label{conj:minor}
For every $\lambda \vdash n$ with $\dim V_\lambda^H(\Sym_n) = 2$ and
$V_\lambda \downarrow \Sym_{n-1} = V_\mu \oplus V_\nu$ (both summands rank-$1$),
the branching minor $\mathcal{M}_\lambda(q)$ defined analogously via
$\Pibs = \Pibs^{(n-1)} \cdot R_n$ with $R_n = (1+T_{n-1}) \cdots (1+T_1)$ satisfies
$\mathcal{M}_\lambda(q) \in \Z_{\geq 0}[q]$ with non-negative palindromic core.
\end{conjecture}

If Conjecture~\ref{conj:minor} holds, the present proof extends inductively
to all rank-$2$ shapes at all $n$ via the recursion
\[ \sigma_2(V_\lambda \mid \Sym_n) = \sigma_1(V_\mu \mid \Sym_{n-1}) \cdot \sigma_1(V_\nu \mid \Sym_{n-1}) \cdot \mathcal{M}_\lambda(q). \]
(Our $\sigma_1$-G1 theorem handles the outer two factors; only the minor remains.)

Our result verifies Conjecture~\ref{conj:minor} for $n = 5$.

\subsection*{Relation to $\sigma_2$-G1 in general}

The proof here reduces $\sigma_2$-G1 at $n = 5$ to $\sigma_1$-G1 at $n = 4$
(already proved) plus positivity of the branching minor (verified by direct
computation). A proof of $\sigma_k$-G1 for all $(k, \lambda, n)$ in
generality would require either:
\begin{itemize}
\item A general non-negativity theorem for $k \times k$ minors of $R_n$ on
appropriate $k$-dimensional branching subspaces (hard), or
\item A categorification of $\sigma_k$ in the KL / Soergel framework
that bypasses determinants entirely (speculative, connected to a-filtration
approach).
\end{itemize}

\subsection*{What this does not close}

The novel degree-$10$ orphan palindrome at $n = 7$, $\lambda = (3,2,1,1)$,
$k = 2$ (memory: \texttt{project\_novel\_atom\_deg10}) involves a shape with
$\dim V_\lambda^H \geq 3$ and so is not covered by the rank-$2$ argument above.
An inductive extension via Conjecture~\ref{conj:minor} would still be
insufficient; rank-$\geq 3$ requires a genuinely new mechanism.

\begin{thebibliography}{9}
\bibitem{KL} D.~Kazhdan and G.~Lusztig, \emph{Representations of Coxeter groups and Hecke algebras}, Invent. Math. 53 (1979), 165--184.
\bibitem{Clio_sigma1} Clio, \emph{Non-negativity of the trace core for the staircase Bott--Samelson operator on Hecke cell modules}, \texttt{2026-04-24-sigma1-G1.pdf}, \url{https://github.com/clio-vega/proofs}.
\bibitem{Hoefsmit} P.~N.~Hoefsmit, \emph{Representations of Hecke algebras of finite groups with BN-pairs of classical type}, Ph.D. thesis, University of British Columbia (1974).
\bibitem{Mathas} A.~Mathas, \emph{Iwahori--Hecke algebras and Schur algebras of the symmetric group}, Univ. Lecture Series 15, AMS (1999).
\end{thebibliography}

\end{document}
